\begin{helper}[Situation 7.3: a Gamma owner absent from the penultimate prefix]
Let $G$ be a finite loopless multigraph with a finite individual edge
type, let $[k]$ be its palette, and fix a distinguished edge $e_0$, an
input coloring, an exact ordered tree-sequence $T$, a generating coloring
$\varphi$, an ETT $\mathcal E$ on these data, a
\noderef{HierarchyBase} $\mathcal B$, a
\noderef{HierarchyCounterexample} $C$, and a
\noderef{HierarchyMinimumCounterexampleContext} $\mathit{ctx}$.
Thus the context supplies properness, exactly-one-uncoloredness, all six
technical assertions at strictly lower positive rungs, fewer-divider
induction at the same rung, rung-local lexicographic minimality, and the
terminal-ladder bound $|\overline\varphi(T_n)|\ge 2n+11$.
Write $n=\mathcal E.\mathit{rounds}$, $H=C.\mathit{hierarchy}$,
$q=H.\mathit{dividerCount}$, and set
\[
 A=\noderef{HierarchyStarredStage}(H,q),\quad
 \ell=\noderef{HierarchyFullTailLength}(C),\quad
 T(y_i)=\noderef{HierarchyFullTailPrefix}(C,i),\quad
 B_{\rm pen}=T(y_{\ell-1}),\quad y_\ell=T.\mathit{vertex}(\mathit{last}).
\]
Let $T_n$ be the terminal ladder tree
$\mathcal E.\mathit{series.data.tree}((\operatorname{Fin.last}n).\mathit{castSucc})$,
and put
\[
 D=\noderef{HierarchyStoredColorsAt}(H,q),\quad
 \Gamma=\bigcup_{\rho\in D}C.\mathit{good}.\Gamma^q_\rho,\quad
 M_0=\bigcup_{v\in V(\mathcal B.\mathit{core})\setminus V(T_n)}
             \noderef{MissingColors}(G,\varphi,v).
\]
Here $V(U)=\noderef{TreeSequenceVertexSet}(U)$ and
$\overline\varphi(U)=\noderef{TreeMissingColors}(\varphi,U)$.
Assume $2\le j<\ell$, where
$j=\noderef{HierarchyPathNumber}(H)$.
Let $\alpha$ be missing on $A$ and at $y_\ell$.
Suppose $\eta\in D$, $\alpha\in C.\mathit{good}.\Gamma^q_\eta$,
and $\eta\notin\overline\varphi(B_{\rm pen})$.
Then there are an original-relative
\noderef{HierarchyReplacementPresentation} $P$ and a color $\beta$
missing under $P.\mathit{coloring}$ both at $y_\ell$ and on the original
exact $X=T(y_{\ell-2})$, such that $\beta$ lies outside the union of
the last Gamma pairs of $P.\mathit{counterexample}$ and the base-extra
missing set of $P$ (using its own base core and terminal ladder).
\end{helper}
\begin{proof}
Apply \noderef{HierarchyCounterexamplePenultimateElementary} to
$\mathit{ctx.proper}$, $\mathit{ctx.oneEdgeUncolored}$,
$\mathit{ctx.lowerRung}$, $\mathit{ctx.fewerDividers}$,
$\mathit{ctx.lexicographicMinimality}$, and
$\mathit{ctx.largeMissingSet}$, with the ambient data in the statement.
It makes $B_{\rm pen}$ elementary under $\varphi$.
The fields of \noderef{TashkinovHierarchy} and the entrywise definition
of \noderef{TreeSequenceInitialSegment} give $A\preccurlyeq T$.
Since $\ell=|T|-|A|\ge3$, the exact prefix $B_{\rm pen}$ has length
$|T|-1$ and contains $A$. The base core is below $A$: at $q=0$
it equals $A$, and at $q>0$ use the core-inside-first field followed
by the nested stages. Also $T_n\preccurlyeq\mathcal B.\mathit{core}$:
in PE this is the left-segment field of the ordered union in
\noderef{HierarchyBase}, and in non-PE the core equals $T_n$.
Consequently $T_n$ and the base core both lie in $B_{\rm pen}$.

By \noderef{GoodHierarchyReservedDisjoint}, every last Gamma pair lies
in $\overline\varphi(A)$ and $\Gamma\cap D=\varnothing$.
Expanding \noderef{HierarchyStoredColorsAt}, one also has
$D=D_n\setminus\overline\varphi(A)$, where
$D_n=\noderef{TashkinovStoredColors}(\mathcal E.\mathit{series.data},n)$.
In particular no stored color is missing on $A$.

The totalized \noderef{HierarchyFirstMissingPosition} of $\eta$ is the
last index of $T$: any earlier actual witness would lie in
$B_{\rm pen}$, contrary to the hypothesis; if no witness exists, the
definition uses the same last index as fallback. The avoidance clause
of \noderef{GoodHierarchy}, together with
\noderef{HierarchySegmentColors}, therefore excludes $\alpha$ from
every exact displayed edge in $T-A$, that is, every edge position
$|A|\le r<|T|$.

Put $X=T(y_{\ell-2})$, let $U$ be the colors on the exact edge interval
$|A|\le r<|X|$, and let $N$ be the colors on the two exact edges in
$T-X$. Then $|N|\le2$ and every color on $T-A$ lies in $U\cup N$.
Apply \noderef{NineN} with the ambient data, $C.\mathit{positive}$,
the base and good hierarchy, the elementary prefix just derived, the
large-set projection, and $i=\ell-2<\ell$. Its available set is
\[
 \mathcal A=(\overline\varphi(X)\setminus M_0)\setminus U .
\]
If $|\mathcal A\setminus(\Gamma\cup D)|\ge5$, removal of $N$ leaves
at least $5-2=3$ colors. Choose $\beta$ in that difference. Then
$\beta$ is missing on $X$, outside $M_0\cup\Gamma\cup D$, and absent
from $U\cup N$.

Otherwise $|\mathcal A\setminus(\Gamma\cup D)|\le4$, so the second
conclusion of NineN supplies an injective family
$(\rho_h)_{h\in\operatorname{Fin}7}$ in $D$, each missing on $X$, with
\[
 Q_h=C.\mathit{good}.\Gamma^q_{\rho_h}\cup\{\rho_h\}
 \quad\hbox{disjoint from }U.
\]
The blocks $Q_h$ are pairwise disjoint. Indeed distinct indices have
distinct owners by injectivity. Their Gamma pairs are disjoint by
\noderef{GoodHierarchy}'s disjoint-pairs field; a Gamma color cannot
equal any stored owner by \noderef{GoodHierarchyReservedDisjoint};
and the two stored singletons are disjoint since their owners differ.
These four membership alternatives prove full block disjointness.
For each block meeting $N$, choose one color in its intersection
with $N$. Disjointness makes these choices injective into $N$.
Thus at most two of the seven blocks meet $N$. Choose a block
meeting none and set $\beta=\rho_h$. The whole chosen block avoids
$U\cup N$, so in particular $\beta$ is absent on $T-A$.
It is missing on $X$, outside $\Gamma$ by reserved disjointness, and
outside $M_0$: the latter set is contained in $\overline\varphi(A)$,
whereas $\beta\in D$ is not missing on $A$.

In both alternatives we have proved
\[
 \beta\in\overline\varphi(X),\quad
 \beta\notin\Gamma\cup M_0,\quad
 \alpha,\beta\hbox{ absent on }T-A,\quad \alpha\ne\beta.
\]
The inequality follows because $\alpha\in\Gamma$ and $\beta\notin\Gamma$.
Both colors have actual missing witnesses in $B_{\rm pen}$.

Apply \noderef{ChangeLemma} to $(\alpha,\beta)$ with the six context
projections, the derived elementary prefix, and these old missing
witnesses. If $\beta$ is missing on $A$, the conditional
ordered-unused premise is vacuous. Otherwise the first $\alpha$
witness is on $A$ and the first $\beta$ witness is strictly after
$A$, so the former precedes the latter in
\noderef{TreeVertexOrder}; absence of $\alpha$ on the full tail verifies
the required edge exclusion through that latter witness. The
avoidance clause for $\alpha\in\Gamma^q_\eta$ puts $\alpha$ in
\noderef{HierarchyStageMissingColors}: at $q>0$ this is missingness on
$A=T_{n,q}$; at $q=0$ it is missingness on $T_n$.
Thus the respective first or second conclusion identifies the
components through the two first missing witnesses.

Those witnesses are in $B_{\rm pen}$ and distinct from $y_\ell$.
\noderef{P2905_T2_ElementaryTerminalComponent}, with the old elementary
prefix and the terminal vertex missing $\alpha$, gives the disjoint
terminal component, a possibly singleton path certificate, and equality
of all missing sets on $B_{\rm pen}$ after its swap $\sigma$.
This handles coincident witnesses as well as different witnesses.
The terminal exchange clause of
\noderef{TwoColorComponentSwapMissingBehavior} makes $\beta$ missing
at $y_\ell$; prefix equality preserves its witness on $X$.
Since neither affected color occurs on the displayed interval $T-A$,
the definition of \noderef{TwoColorComponentSwap} fixes every exact
edge value on that entire interval. No claim about arbitrary edges
inside $A$ follows from the missing-set equality.
The entire missing set on $B_{\rm pen}$ is now fixed. Therefore every
totalized \noderef{HierarchyFirstMissingPosition} on $T$ is fixed,
in particular that of every $\rho\in D$. Indeed, if a color has a
witness before the last vertex, equality at every preceding vertex
preserves exactly its nonempty set of early candidates and its least
candidate. If it has no early witness, neither coloring has one,
and both totalized positions are the last index, whether or not a
terminal witness exists. This checks the fallback as well as the
actual-witness case.

All premises of \noderef{StableChangeCasesOneTwoFullTail} have now
been verified: distinct colors missing on the old penultimate prefix,
the terminal path and swap, the Cases 1/2 alternative with distinguished
color $\alpha$, disjointness from its first-witness component, fixed
values on every displayed tail edge, pointwise penultimate missing-set
equality, and equality of every stored first position. Apply that
principal and retain its full-tree corresponding maximum ETT, rebuilt
base, good hierarchy, and same-good witness. FullTail is essential
here: the owner $\eta$ has no penultimate witness, so its cutoff need
not be below any strict split prefix.
We describe the identity presentation used below. For every slot of
the certified \noderef{TashkinovSeries}, properness and exact
one-uncoloredness follow by cases on its
\noderef{TashkinovSeriesStep}: the initial case uses the two coloring
clauses of its \noderef{KTriple}; SE uses
\noderef{MaximumDefectiveChoice}'s selected-coloring clauses; PE uses
the explicit current-coloring clauses; RE copies the previous coloring.
Choose the same input, series, and ETT for the reflexive
\noderef{CorrespondingExtendedTashkinovTree}; use reflexive rung,
presentation, ladder, stored-set, connecting-edge, and kind equalities,
the original generating-coloring equality, the just verified slotwise
coloring laws, and reflexive \noderef{StableColoring} at each positive
slot. Define $P_0$ by coloring $\varphi$, that certificate, maximum
$C.\mathit{maximum}$, base $\mathcal B$, counterexample $C$, reflexive
\noderef{SameGoodHierarchy}, and the supplied context $\mathit{ctx}$.
These are all seven fields of
\noderef{HierarchyReplacementPresentation}.
Retain the provider's full corresponding certificate $\mathcal C$,
its maximum witness, base $\mathcal B'$, hierarchy $K$, good witness
$g$, and $\noderef{SameGoodHierarchy}(C.\mathit{good},g)$.
Construct a \noderef{HierarchyCounterexample} $C'$ on $\mathcal B'$:
its positive-rung field is $C.\mathit{positive}$ transported by
$\mathcal C.\mathit{rungAgreement}$, its hierarchy and goodness are
$K,g$, and its maximum is the provider's maximum witness.
The two stated missing witnesses, one in the strict prefix
$B_{\rm pen}$ and one at $y_\ell$, are distinct by injectivity of
\noderef{TreeSequence}'s vertex display. Their common missing color
contradicts the disjointness clause of
\noderef{ColoredVertexSetElementary}, proving the non-elementarity
field. Apply \noderef{HierarchyReplacementPresentationTrans} to $P_0$,
$\sigma$, $\mathcal C$, $\mathcal B'$, $C'$, and the supplied same-good
witness. It gives an original-relative presentation $P$ with coloring
exactly $\sigma$; in particular its minimum context is supplied by the
principal, not assumed.

For any such $P$, its own same-good field identifies its base core,
last stored set, and every last Gamma pair with the originals after the
divider-count cast. The rung and ladder agreement in its corresponding
certificate identifies its actual terminal ladder tree with $T_n$:
use the last positive slot, which exists since $C.\mathit{positive}$
gives $n>0$, and its series-presentation equality. Missing-set equality
on the old base core and $T_n$ then identifies its base-extra missing
set with $M_0$. Thus the final forbidden set in the statement is the
original $\Gamma\cup M_0$.

Here use $\beta$'s retained witness on $X\subseteq B_{\rm pen}$ to
construct the non-elementarity field. In both the free-color and the
stored-owner choices $\beta\notin\Gamma\cup M_0$, as proved above.
The final forbidden-set identification therefore gives exactly the
asserted exclusion for $P,\beta$, and its two missing witnesses
give the rest of the conclusion.
\end{proof}
